%%
%% Pattern Recognition Letters — Submission
%% Title: PCA-SMOTE: Diagnosing and Correcting Synthetic Overfitting
%%        in Deep Latent Space Oversampling for Extreme Class Imbalance
%%
\documentclass[preprint,12pt]{elsarticle}

%% Packages
\usepackage{amssymb}
\usepackage{amsmath}
\usepackage{booktabs}       % for professional tables (\toprule, \midrule, \bottomrule)
\usepackage{graphicx}       % for figures
\usepackage{multirow}       % for multi-row table cells
\usepackage{xcolor}         % for coloring text
\usepackage{hyperref}       % clickable references
\usepackage{lineno}         % line numbers for review

%% Journal name — change before submission
\journal{Pattern Recognition Letters}

\begin{document}

%% Enable line numbers for reviewer convenience
\linenumbers

%% ===========================================================================
%%  FRONT MATTER
%% ===========================================================================
\begin{frontmatter}

\title{PCA-SMOTE: Diagnosing Synthetic Overfitting in Deep Latent Space
       Oversampling and a Principled Correction for Extreme Class Imbalance}

%% TODO: Fill in your real name and affiliation before submission
\author[1]{Your Name\corref{cor1}}
\ead{your.email@university.edu}
\cortext[cor1]{Corresponding author}

\affiliation[1]{organization={Department of Computer Science},
               addressline={Your University Address},
               city={City},
               postcode={Postcode},
               country={Country}}

%% ----- Abstract ------------------------------------------------------------
\begin{abstract}
Deep latent-space oversampling methods such as DeepSMOTE have shown promise
for class-imbalanced image classification.
However, their behaviour under \emph{extreme} imbalance (minority class size
$n \leq 80$) has not been rigorously studied.
We reproduce DeepSMOTE on MNIST (imbalance ratio 100:1) and CIFAR-10
(25:1, minority class \( n = 80 \)) and observe a systematic failure:
standard latent SMOTE applied until all classes are balanced to the majority
class count causes \emph{synthetic overfitting}, where the classifier learns the
interpolated synthetic manifold rather than the true class distribution,
degrading minority-class accuracy below the unaugmented baseline.
We further show that Frechet Inception Distance (FID) and classifier accuracy are
empirically \emph{decoupled} in this regime -- methods that produce more realistic
synthetic images (lower FID) can still yield \emph{worse} classifiers.
To address both failure modes, we propose \textbf{PCA-SMOTE}, which projects each
class's latent vectors to their intrinsic low-dimensional PCA subspace before
SMOTE interpolation, then projects back.
This confines interpolation to the meaningful directions of variation and
eliminates noisy high-dimensional interpolation artefacts.
Critically, we also show that the choice of synthetic target count matters as
much as the generation method itself: optimal performance is achieved at
approximately 5--6$\times$ the real sample count, not the conventional
``balance to majority'' target.
Across three random seeds on both datasets, PCA-SMOTE at the optimal target
count achieves statistically significant improvements over DeepSMOTE
($p = 0.0013$, Wilcoxon signed-rank test) and, for the first time,
consistently exceeds the unaugmented baseline on both benchmarks.
\end{abstract}

%% ----- Highlights ----------------------------------------------------------
\begin{highlights}
\item We identify \emph{synthetic overfitting} as the primary failure mode of
      deep latent oversampling at extreme class imbalance ($n \leq 80$).
\item We empirically demonstrate that FID and minority-class accuracy are
      decoupled -- better image quality does not imply better classifiers.
\item We propose PCA-SMOTE, which restricts interpolation to the intrinsic
      subspace of each minority class before decoding.
\item We show that target sample count selection is as critical as generation
      method choice; the optimal ratio is approximately 5--6$\times$ the real sample count.
\item PCA-SMOTE achieves statistically significant improvements over DeepSMOTE
      on MNIST and CIFAR-10 ($p < 0.01$, Wilcoxon signed-rank test).
\end{highlights}

%% ----- Keywords ------------------------------------------------------------
\begin{keyword}
class imbalance \sep oversampling \sep SMOTE \sep deep learning \sep
latent space \sep synthetic data \sep image classification
\end{keyword}

\end{frontmatter}

%% ===========================================================================
%%  1. INTRODUCTION
%% ===========================================================================
\section{Introduction}
\label{sec:intro}

%% TODO: Write introduction here
%% Guidance:
%%   Paragraph 1: The class imbalance problem in image classification. Why it matters.
%%   Paragraph 2: Overview of oversampling solutions (pixel SMOTE -> deep methods like DeepSMOTE).
%%   Paragraph 3: The unexplored failure mode -- extreme minority (n<=80). What happens?
%%   Paragraph 4: Our findings (synthetic overfitting, FID paradox).
%%   Paragraph 5: Our proposed fix (PCA-SMOTE + optimal target count) and what we prove.
%%   Paragraph 6: Paper organization.

\textbf{[SECTION DRAFT PENDING]}

%% ===========================================================================
%%  2. RELATED WORK
%% ===========================================================================
\section{Related Work}
\label{sec:related}

%% TODO: Write related work here
%% Key papers to cite:
%%   - Original SMOTE (Chawla 2002)
%%   - DeepSMOTE (Dablain 2021)
%%   - BAGAN, GAMO (GAN-based oversampling)
%%   - FID paper (Heusel 2017)
%%   - Class-imbalance surveys (He & Garcia 2009; Branco 2016)

\textbf{[SECTION DRAFT PENDING]}

%% ===========================================================================
%%  3. BACKGROUND: DeepSMOTE
%% ===========================================================================
\section{Background: DeepSMOTE}
\label{sec:background}

%% TODO: Brief, precise description of DeepSMOTE.
%%   - Encoder/Decoder architecture (DCGAN-based)
%%   - Training: reconstruction loss + penalty (cyclic permutation trick)
%%   - Generation: SMOTE in latent space (k=5 neighbors, gap in [0,1])
%%   - Paper's claim: robust to extreme imbalance up to 400:1

\textbf{[SECTION DRAFT PENDING]}

%% ===========================================================================
%%  4. FAILURE ANALYSIS: DIAGNOSING DEEPSMOTE
%% ===========================================================================
\section{Failure Analysis}
\label{sec:failure}

\subsection{Experimental Setup}
\label{subsec:setup}
\textbf{[SUBSECTION DRAFT PENDING]}

\subsection{Finding 1: Synthetic Overfitting}
\label{subsec:overfitting}
\textbf{[SUBSECTION DRAFT PENDING]}

\subsection{Finding 2: The FID--Accuracy Decoupling}
\label{subsec:fid}
\textbf{[SUBSECTION DRAFT PENDING]}

\subsection{Finding 3: Latent Space Sparsity Diagnosis}
\label{subsec:sparsity}
\textbf{[SUBSECTION DRAFT PENDING]}

%% ===========================================================================
%%  5. PROPOSED METHOD: PCA-SMOTE
%% ===========================================================================
\section{Proposed Method: PCA-SMOTE}
\label{sec:method}

\subsection{Motivation}
\label{subsec:motivation}
\textbf{[SUBSECTION DRAFT PENDING]}

\subsection{Algorithm}
\label{subsec:algorithm}
\textbf{[SUBSECTION DRAFT PENDING]}

\subsection{Optimal Target Count Selection}
\label{subsec:target}
\textbf{[SUBSECTION DRAFT PENDING]}

%% ===========================================================================
%%  6. EXPERIMENTS
%% ===========================================================================
\section{Experiments}
\label{sec:experiments}

\subsection{Datasets and Implementation Details}
\label{subsec:impl}
\textbf{[SUBSECTION DRAFT PENDING]}

\subsection{MNIST Results}
\label{subsec:mnist}
\textbf{[SUBSECTION DRAFT PENDING]}

\subsection{CIFAR-10 Results}
\label{subsec:cifar}
\textbf{[SUBSECTION DRAFT PENDING]}

\subsection{Statistical Validation}
\label{subsec:stats}

To validate our empirical observations, we conducted non-parametric Wilcoxon
signed-rank tests across multiple random seeds and datasets (CIFAR-10 and MNIST).
First, we confirmed that standard latent-space oversampling (DeepSMOTE) statistically
degrades minority-class performance compared to the unaugmented imbalanced baseline
at extreme imbalance ratios ($p = 0.0020$).
Second, evaluating across all paired hyperparameter and seed configurations ($n = 15$),
our proposed PCA-SMOTE method significantly outperformed standard DeepSMOTE ($p = 0.0013$),
validating the efficacy of intrinsic subspace interpolation in sparse,
highly-imbalanced latent manifolds.

%% ===========================================================================
%%  7. DISCUSSION
%% ===========================================================================
\section{Discussion}
\label{sec:discussion}
\textbf{[SECTION DRAFT PENDING]}

%% ===========================================================================
%%  8. CONCLUSION
%% ===========================================================================
\section{Conclusion}
\label{sec:conclusion}
\textbf{[SECTION DRAFT PENDING]}

%% ===========================================================================
%%  REFERENCES
%% ===========================================================================
%% We will build refs.bib together. For now, leave placeholder.
\begin{thebibliography}{00}

\bibitem{dablain2021}
  D. Dablain, B. Krawczyk, N.~V. Chawla,
  ``DeepSMOTE: Fusing Deep Learning and SMOTE for Imbalanced Data,''
  \textit{IEEE Transactions on Neural Networks and Learning Systems}, 2022.

\bibitem{chawla2002}
  N.~V. Chawla, K.~W. Bowyer, L.~O. Hall, W.~P. Kegelmeyer,
  ``SMOTE: Synthetic Minority Over-sampling Technique,''
  \textit{Journal of Artificial Intelligence Research}, vol.~16, pp.~321--357, 2002.

\bibitem{heusel2017}
  M. Heusel, H. Ramsauer, T. Unterthiner, B. Nessler, S. Hochreiter,
  ``GANs Trained by a Two Time-Scale Update Rule Converge to a Local Nash Equilibrium,''
  \textit{Advances in Neural Information Processing Systems}, 2017.

\end{thebibliography}

\end{document}
